\section{Introduction}

Single-cell RNA sequencing (scRNA-seq) has changed the scale at which cellular heterogeneity can be studied. Droplet technologies now support measurements from hundreds of thousands of cells, and coordinated projects such as the Human Cell Atlas and Tabula Sapiens have made multi-organ reference maps a routine part of biomedical research \citep{Macosko2015,Zheng2017,Regev2017,TabulaSapiens2022,Tanay2017}. These atlases are valuable only when transcriptional profiles can be assigned to biologically meaningful cell identities. Cell-type annotation is therefore not a cosmetic post-processing step: it determines the units used for differential abundance, state-specific expression, trajectory analysis, cell--cell communication, and association with clinical phenotypes. Manual marker-based annotation remains widely used, but it is labour intensive, sensitive to the available expertise and marker list, and difficult to reproduce at atlas scale. General best-practice and reporting frameworks consequently treat annotation provenance and reproducible preprocessing as essential parts of a single-cell analysis \citep{LueckenTheis2019,Lahnemann2020,Fullgrabe2020}.

Automated annotation methods have approached this problem through several statistical traditions. SingleR and scmap project cells onto labelled references; scPred uses a supervised support-vector framework; CellTypist uses regularized logistic regression; Azimuth combines reference mapping with harmonized single-cell resources; and latent-variable models such as scVI learn lower-dimensional representations that can support downstream annotation \citep{Aran2019,Kiselev2018,AlquicirayHernandez2019,DominguezConde2022,Hao2021,Lopez2018}. Cell BLAST, scNym, Harmony and scArches further illustrate the range from learned search spaces and semi-supervised classification to integration and transfer learning \citep{CellBlast2020,ScNym2021,Korsunsky2019,Lotfollahi2022}. Broad benchmarks have repeatedly shown that relatively simple supervised models can be difficult to outperform when training and test data share protocols and labels \citep{Abdelaal2019,Pasquini2021}. Their success does not remove the need for transferable representations, because curated labels can be scarce, datasets can differ by donor or platform, and cell states may not align cleanly between references. It does, however, establish that new representation-learning methods require carefully matched baselines.

Single-cell foundation models extend this representation-learning idea by pretraining transformer-like architectures on millions of transcriptomes. scBERT, Geneformer, scGPT, scFoundation, Universal Cell Embeddings and CellFM use different tokenization schemes and pretraining objectives, but share the aim of learning reusable cell or gene representations \citep{Yang2022,Theodoris2023,Cui2024,Hao2024,Rosen2023,CellFM2025,Baek2025}. scGPT is a prominent example. Its whole-human checkpoint was trained on a large CELLxGENE-derived corpus and exposes an interface for obtaining fixed-dimensional cell embeddings that can be used by a separate task-specific learner \citep{Cui2024,CZICELLxGENE2023}. Pretraining could be most useful when target labels are limited, when biological variation must be transferred between studies, or when the downstream task benefits from structure learned outside the target cohort. Yet pretraining alone does not produce an annotation probability. A frozen encoder must be coupled to a classifier, and the resulting performance and calibration belong to that complete pipeline.

This distinction is central to fair evaluation. If scGPT embeddings are paired only with a linear classifier while expression features are paired with a nonlinear learner, the comparison confounds representation and classifier capacity. Conversely, a classical model given a larger or differently selected gene set may benefit from an information advantage. Feature selection performed before a train--test split can leak information from the test partition even when class labels are not used. Sparse matrices create a further implementation hazard for some tree libraries because omitted biological zeros can be interpreted as missing values. A controlled benchmark should therefore match the genes available to both representations, apply the same classifier families, learn all filters within the training partition, preserve biological zeros, and name each representation--classifier combination explicitly. These safeguards also align with broader calls for transparent reporting of single-cell experiments and computational transformations \citep{Fullgrabe2020}.

Generalization also depends on the unit of splitting. Cells sampled from one donor share genetics, exposures, processing, and technical effects; assigning cells from the same donor to both training and test sets can yield an optimistic estimate of performance on unseen individuals. The problem is not unique to scRNA-seq, but is amplified by the large ratio of cells to biological replicates. Cell-level bootstrap intervals likewise treat correlated cells as independent and can appear much narrower than uncertainty obtained by resampling donors. Benchmarks of single-cell integration and foundation-model embeddings have shown that batch and donor structure can dominate apparent biological organization \citep{Luecken2022,Kedzierska2025}. More generally, treating cells as independent replicates can produce misleading precision when biological replication occurs at sample or donor level \citep{Squair2021}. Donor-held-out evaluation is therefore a necessary complement to large within-atlas cell-level comparisons, although it still does not substitute for external-atlas transfer when protocols and annotation systems differ.

Probability quality is a second, distinct requirement. Accuracy and Macro F1 describe class discrimination, but many annotation workflows also use confidence scores to reject uncertain cells, prioritize manual review, or weight downstream analyses. Calibration asks whether predicted probabilities correspond to empirical frequencies. Expected calibration error (ECE) is intuitive but depends on bin definitions and is not a proper scoring rule \citep{Guo2017,Vaicenavicius2019,SilvaFilho2023}. Adaptive and classwise ECE expose different aspects of multiclass calibration, whereas the Brier score and negative log-likelihood evaluate the full probability vector and reward honest probabilistic prediction \citep{Gneiting2007,Nixon2019}. Reliability diagrams are most interpretable when bin counts and uncertainty are shown. Post-hoc approaches range from a single temperature to multiclass calibrators, but require calibration observations that remain separate from evaluation data \citep{Kull2019,VanCalster2019}. Temperature scaling provides a simple test of whether probability sharpness, rather than class ranking, accounts for observed miscalibration \citep{Guo2017}.

Independent evaluations have begun to question whether large single-cell models consistently improve over simpler representations. Zero-shot studies found that scGPT and Geneformer did not uniformly outperform highly variable genes, Harmony, or scVI for cell-type structure and batch integration, and that results depended strongly on dataset and task \citep{Kedzierska2025}. Other benchmarking frameworks have emphasized inconsistent software interfaces, hyperparameter sensitivity, and the need to distinguish zero-shot, frozen-probe, and fine-tuned settings \citep{Boiarsky2023,BioLLM2025,ScEval2026}. These observations do not imply that pretraining lacks value. They instead reinforce a general principle: the benefit of a foundation model must be demonstrated for a specified representation, downstream learner, label regime, and biological split.

Our original submission addressed accuracy and calibration across six large atlases but used a single cell-level split, test-informed feature selection, a custom scGPT wrapper, and an incomplete calibration assessment. A full executable audit during revision showed that the wrapper did not reproduce the released scGPT embedding recipe and that sparse XGBoost inputs could treat zero counts as missing. We therefore superseded the earlier values and reran the benchmark from raw counts using the official scGPT 0.2.5 embedding helper, training-only Seurat v3 feature selection, exact vocabulary matching, dense XGBoost inputs, and a factorial comparison of two representations by three classifiers. This correction narrows the claims but makes the numerical record reproducible.

The revised study addresses four questions. First, how do official frozen scGPT embeddings compare with normalized expression of the same genes when both are paired with Logistic Regression, Random Forest, and XGBoost? Second, how do conclusions change when donors rather than cells define five independent train--test partitions? Third, does reducing the number of labelled training cells reveal an advantage for the pretrained representation? Fourth, how stable are calibration conclusions across complementary metrics and after temperature scaling fitted without access to test donors? We answer these questions across six public atlases containing 1,405,933 cells, with a focused donor-held-out analysis in 199 Indonesian donors. The results support strong classical baselines for the tested supervised settings while defining, rather than dismissing, the transfer regimes that remain open.
